% Testdokument für Spike S2 (pdf.js-Textebene für Markierungsanker). Jede Zeile ist bewusst so gesetzt,
% dass die Testmarkierungen in marks.json auf sie zeigen können. Zeilennummern nicht verschieben!
\documentclass[11pt,a4paper]{article}
\usepackage[T1]{fontenc}
\usepackage[utf8]{inputenc}
\usepackage{lmodern}
\usepackage[ngerman]{babel}
\usepackage{microtype}
\usepackage{amsmath,amssymb}
\usepackage{booktabs}
\usepackage{csquotes}
\usepackage[hidelinks]{hyperref}
\begin{document}

\section{Fließtext und Ligaturen}

Die offizielle Auflage der Zeitschrift beschreibt die Effizienz der Schifffahrt und die Pflichten der Besatzung.
Ein flacher Grafikfilter findet auffällige Stellen im Fluss der Wörter, ohne dass jemand eingreifen muss.
Weiterhin gilt: Fünf schöne Bücher über Größe, Maß und Übermaß liegen auf dem Tisch neben der Kaffeetasse.

Dieser zweite Absatz ist absichtlich lang, damit der Zeilenumbruch mitten im Satz geschieht und eine Markierung
über mehrere Zeilen reichen muss. Die Wörter am Ende einer Zeile und am Anfang der nächsten Zeile sind dabei
besonders interessant, weil die Textebene von pdf.js dort Zeilenenden einfügt und die Lage der Zeichen von der
Schriftart abhängt, die der Browser als Ersatz verwendet. Am Schluss steht ein kurzer Satz.

\section{Formeln}

Für $x \in \mathbb{R}$ gilt $f(x) = x^2 + 2x + 1$, also ist die Funktion niemals negativ und hat genau eine Nullstelle.
Die abgesetzte Gleichung
\begin{equation}
    \int_0^1 x^2 \, dx = \frac{1}{3}
\end{equation}
zeigt ein Integral, danach geht der Text ganz normal weiter und endet mit einem Punkt.

\section{Silbentrennung}

\begin{minipage}{0.42\textwidth}
Die Donaudampfschifffahrtsgesellschaft beschäftigt viele Rechtsschutzversicherungsgesellschaften und
Haftpflichtversicherungen, deren Vertragsbedingungen ausgesprochen unübersichtlich formuliert sind.
\end{minipage}

\section{Fußnote und Auszeichnung}

Hier steht ein Satz mit einer Fußnote\footnote{Die Fußnote enthält zusätzlichen Text über Markierungen im PDF.} und geht danach weiter.
Ein \emph{kursives Wort} und \textbf{fette Schrift} sowie \texttt{Schreibmaschine} stehen mitten im Satz.
Zitate wie \enquote{dieses Beispiel} und Zahlen wie 3,14 oder 2026 sowie z.\,B. Abkürzungen kommen ebenfalls vor.

\section{Tabelle}

\begin{tabular}{@{}llr@{}}
    \toprule
    Werkzeug & Aufgabe & Anteil \\
    \midrule
    Editor & Schreiben des Textes & 60 \\
    Compiler & Erzeugen der Datei & 40 \\
    \bottomrule
\end{tabular}

\section{Aufzählung}

\begin{itemize}
    \item Der erste Punkt der Liste
    \item Der zweite Punkt mit etwas mehr Text
\end{itemize}

\end{document}
